\section{A high-depth prime at the maximal actual rank}
\label{hr-section}

The net sign of a rank packet differs from the occurrence of a single
high-depth prime in that packet. We give actual homogeneous profiles
with small $\lambda$ for which the latter event occupies the maximal
rank. The construction does not determine the net packet sign.

Write $\mathcal O=\mathbb Z[\zeta]$, $\zeta^2-\zeta+1=0$,
$P(A+B\zeta)=AB(A+B)$ and $N(A+B\zeta)=A^2+AB+B^2$.
For an unramified primitive root $w$ with $N(w)>1$, and $p>3$ prime
to $N(w)$, let
$d_p=\operatorname{ord}((w/\bar w)^3\bmod p)$ and
$s_p=v_p(P(w^{d_p}))$. At a cutoff $Y$, distinguish
\begin{align*}
 E_{\rm net}(g,Y)&=\left\{d\mid g:
    \sum_{\substack{p>Y\\d_p=d}}(s_p-3)\log p>0\right\},\\
 E_{\rm hit}(g,Y)&=\{d\mid g:\text{ some }p>Y
                            \text{ has }d_p=d,\ s_p\ge4\}.
\end{align*}
These sums and sets range only over $p>3$ with $p\nmid N(w)$,
as in the rank definition, even if an initial cutoff is below three.
One always has $E_{\rm net}\subseteq E_{\rm hit}$, but the reverse
inclusion need not follow. Our counterfamily addresses only automatic
totient sparsity of the larger set.

\begin{lemma}[An elementary prime with the required order]\label{hr-prime}
For each prime $\ell\ge5$ there is a prime $p\equiv1\pmod{6\ell}$.
\end{lemma}
\begin{proof}
Put $n=6\ell$ and take a prime divisor $p$ of $\Phi_n(n)$.
Its complex product has absolute value at least $(n-1)^{\varphi(n)}>1$.
The constant term is one, so no prime dividing $n$ divides this value.
In particular $p\nmid n$. Modulo $p$, the factorization
$X^n-1=\prod_{d\mid n}\Phi_d(X)$ is squarefree. The residue $n$
is a root of $\Phi_n$ and of no factor at a proper divisor of $n$.
It therefore has multiplicative order $n$, giving $n\mid p-1$.
\end{proof}
Only this elementary existence is needed, with no least-prime estimate.

\begin{theorem}[Prescribed first depth at a maximal rank]\label{hr-depth}
Fix a prime $\ell\ge5$, a prime $p\equiv1\pmod{6\ell}$ and $s\ge4$.
There is a positive integer $a\equiv0\pmod3$ such that $w=a+\zeta$
satisfies
\[
 \gcd(N(w),3p)=1,\qquad d_p=\ell,\qquad s_p=s,
 \qquad p^s\le a<3p^{s+1}.
\]
\end{theorem}
\begin{proof}
In $\mathbb F_p$ choose $z$ of order six and $R$ of order $\ell$,
and put $\bar z=1-z$. Set
\[
 a_0=\frac{R\bar z-z}{1-R}.
\]
Then $a_0+\bar z=(\bar z-z)/(1-R)\ne0$ and
$a_0+z=R(a_0+\bar z)\ne0$. Their cubed ratio has exact order
$\ell$, and their product is nonzero. For the integer polynomial
$G_\ell(T)=P((T+\zeta)^\ell)$ the cubic identity gives
\[
 3(z-\bar z)G_\ell(T)=(T+z)^{3\ell}-(T+\bar z)^{3\ell}
 \quad\text{over }\mathbb F_p.
\]
At $a_0$ the right side vanishes. If $R(T)=(T+z)/(T+\bar z)$,
its derivative there is the derivative of
$(T+\bar z)^{3\ell}(R(T)^{3\ell}-1)$, which is nonzero since
$p\nmid3\ell$ and
$R'(a_0)=(\bar z-z)/(a_0+\bar z)^2\ne0$.
Thus $a_0$ is a simple root of $G_\ell$.

Lift it to the unique root $b$ modulo $p^s$, with $0\le b<p^s$.
Of the $p$ residues $b+jp^s$ modulo $p^{s+1}$, exactly one lifts
to a root. Choose a nonlifting $j\in\{1,2\}$. Then
$A=b+jp^s$ has exact depth $s$ and $p^s\le A<3p^s$.
Choose $k\in\{0,1,2\}$ by CRT so that
$a=A+kp^{s+1}\equiv0\pmod3$. Since $p>3$, this preserves the
depth and gives the stated upper bound. The norm is a unit modulo
$p$ by the first calculation, and it is one modulo three.
\end{proof}

\begin{theorem}[Failure of automatic coarse rank sparsity]\label{hr-sparsity}
There is a sequence of actual primitive homogeneous profiles $w^g$,
with unit residual and $g\to\infty$, for which $\lambda=1/g\to0$
and, at the actual cutoff $Y=g^{1/6}$,
\[
 \frac1g\sum_{d\in E_{\rm hit}(g,Y)}\varphi(d)
       \ge\frac{g-1}{g}\longrightarrow1.
\]
Thus this coarse exceptional-rank mass is not automatically $o(g)$.
\end{theorem}
\begin{proof}
Take $g=\ell$ through the primes at least five and apply the previous
two results, for example with $s=4$. The root $w=a+\zeta$ is
primitive and has norm prime to three. Oriented factorization has
no inert prime and no pair of conjugate split factors, so every
$w^g$ is primitive. A root-of-unity quotient $w/\bar w$ would make
their coprime ideals equal and force $w$ to be a unit. Its norm
exceeds one, so no power has zero boundary. A unit rotation gives
an actual positive primitive triple and preserves the absolute
boundary and the norm. The profile has unit residual.

The selected prime satisfies $p\ge6g+1>g^{1/6}$, $d_p=g$ and
$s_p=4$. Therefore $g\in E_{\rm hit}$, while $\varphi(g)=g-1$.
\end{proof}

This construction quantifies the loss in marking a whole rank after
seeing one deep prime. For its positive rotation with $t=\log c$,
the size choice in Theorem~\ref{hr-depth} gives
\[
 N(w)\ge p^{2s},\qquad t\ge\frac g2\log N(w)\ge gs\log p,
 \qquad \frac{(s-3)\log p}{t}\le\frac{s-3}{gs}\le\frac1g.
\]
An arbitrarily large prescribed first depth can therefore occur at
the maximal totient rank while its own signed cost is sublinear in
height. Other primes in the same rank may supply negative credits;
their full packet sign and size have not been determined here.

In particular this is not a counterexample to a condition on
$E_{\rm net}$, a rankwise saving hypothesis, or the global signed
tail. The selected prime moves beyond the actual cutoff, but every
other retained contribution must still be accounted for. Exact
integer replay checks fourteen profiles with varied ranks and depths,
including primitivity, positive unit rotations, exact first depth and
all displayed size premises. It does not factor the whole boundary
or claim a net packet sign. No ABC proof or disproof follows.
